% titlepage
\begin{titlepage}
	\begin{figure}[H]
		\centering
		\includegraphics[width = 7cm, keepaspectratio=1]{pics/logo.pdf}\\[12pt]
		{\sc Oddelek za fiziko}\\[4cm]
	\end{figure}
	\begin{center}
		\large{Seminar I$_b$ -- 1. letnik, II stopnja}\\[0.5cm]
		\LARGE\textbf{Supersimetrija v nerelativisti\v cni kvantni mehaniki}\\[1.0cm]

		\vspace{0.0cm}

		\begin{minipage}{0.4\textwidth}\small
			\begin{flushleft}
			\textsc{Avtor:}\\[0.2cm]
			Jo\v ze Zobec, dipl. fiz. (UN)
			\end{flushleft}
		\end{minipage}
		\begin{minipage}{0.4\textwidth}\small
			\begin{flushright}
				\textsc{Mentor:}\\[0.2cm]
				Prof. Dr. Svjetlana Fajfer,\\[0.5cm]
				\textsc{Somentor:}\\[0.1cm]
				Prof. Dr. Toma\v z Prosen
			\end{flushright}
		\end{minipage}
	\end{center}

	\vspace{3.5cm}

	\begin{abstract}
		Supersimetrija je podro\v cje, ki je nastalo za potrebe velike teorije poenotenja,
		vendar pa lahko najde rabo tudi drugod v fiziki. V tem seminarju bom pokazal kako
		supersimetri\v cno matemati\v cno ogrodje uporabimo v \v solskih kvantno-mehanskih
		problemih, v teoriji izospektralnih Hamiltonianov in v teoriji perturbacij.
	\end{abstract}
	
	\vfill

	\centering{\footnotesize Ljubljana, \today}
\end{titlepage}

